\chapter{Quality, Security and Delivery}
\label{ch:quality}

\section{Testing strategy}

\begin{table}[h]
\centering
\small
\begin{tabularx}{\textwidth}{l X l}
\toprule
\textbf{Level} & \textbf{What it covers} & \textbf{Gate} \\
\midrule
Unit & Pure logic: rope operations, solver, graph queries, policy decisions & Every commit \\
Property & Invariants under random input: edits, undo, selections, graph consistency & Every commit \\
Golden & Compression output, SCP records, context assembly for fixture repositories & Every commit \\
Integration & Service and interface across the transport, with real stores & Every commit \\
Interface & Surfaces driven by projections, asserting emitted intents, no window & Every commit \\
End to end & A real mission on a fixture repository with a scripted agent & Every merge \\
Live & A real mission with a real model provider on a sample project & Before release \\
Soak & A long mission with induced failures and a killed process & Before release \\
\bottomrule
\end{tabularx}
\caption{Test levels and when each gate runs.}
\end{table}

Two rules matter more than coverage numbers. Every failure path has a test, because failure
handling is the product. And every test that touches an agent uses a scripted agent by
default, so the suite is deterministic and free; real providers appear only in the live gate.

\section{Performance budgets}

Budgets are enforced by an automated benchmark that fails the build on regression beyond the
stated tolerance.

\begin{table}[h]
\centering
\small
\begin{tabularx}{\textwidth}{X l l}
\toprule
\textbf{Operation} & \textbf{Budget} & \textbf{Measured on} \\
\midrule
Cold start to an interactive window & Under one second & Reference machine, warm cache \\
Open a file of one hundred thousand lines & Under two hundred milliseconds & Reference machine \\
Keystroke to rendered frame & Under sixteen milliseconds at the ninety ninth percentile & Reference machine \\
Incremental reparse after an edit & Under five milliseconds for a typical file & Reference machine \\
Structural impact query & Under fifty milliseconds on a repository of ten thousand files & Warm index \\
Context assembly for a ticket & Under one second excluding model latency & Warm index \\
Snapshot creation for verification & Under two seconds with a warm worktree pool & Reference machine \\
Memory at rest with a large repository indexed & Under one gigabyte & Reference machine \\
\bottomrule
\end{tabularx}
\caption{Performance budgets. A regression beyond tolerance fails the build.}
\end{table}

\section{Security model}

\subsection{Trust boundaries}

\begin{figure}[h]
\centering
\begin{tikzpicture}[
  zone/.style={draw=xyraline,fill=xyrapanel,rounded corners=3pt,font=\scriptsize,align=center,minimum height=11mm},
  danger/.style={draw=xyrared,line width=0.9pt,fill=white,rounded corners=3pt,font=\scriptsize,align=center,minimum height=11mm},
  arr/.style={-{Stealth[length=1.8mm]},xyragrey,font=\scriptsize}
]
\node[zone,minimum width=32mm] (dev) {Developer\\trusted};
\node[zone,minimum width=32mm,right=10mm of dev] (core) {Core service\\trusted};
\node[danger,minimum width=32mm,right=10mm of core] (agent) {Agent output\\untrusted};
\node[danger,minimum width=32mm,below=9mm of agent] (ext) {Extensions\\untrusted};
\node[danger,minimum width=32mm,below=9mm of core] (repo) {Repository content\\untrusted};

\draw[arr] (dev) -- (core);
\draw[arr] (core) -- node[above] {brief} (agent);
\draw[arr] (agent) -- node[below,xshift=6mm] {diff, tool calls} (core);
\draw[arr] (ext) -- node[right] {capability calls} (core);
\draw[arr] (repo) -- node[right] {content} (core);
\end{tikzpicture}
\caption{Trust boundaries. Everything a model or a repository produces is untrusted input.}
\end{figure}

\subsection{Controls}

\begin{longtable}{@{}p{0.30\textwidth} p{0.62\textwidth}@{}}
\toprule
\textbf{Risk} & \textbf{Control} \\
\midrule
\endhead
Prompt injection from repository content & Repository text is data, never instruction. Tool calls originating from content patterns are rejected, and the Context Inspector shows provenance for every item \\
Secret exfiltration through context & Redaction runs before any content leaves the machine, covering keys, tokens, environment files and private keys, and the redaction is recorded \\
Malicious extension & Sandboxed execution with an explicit capability list, no ambient filesystem or network, and a trap disables the extension for the session \\
Destructive agent action & Write access is a property of the role, snapshots are used for execution, and every mutation is a reversible transaction with a commit boundary \\
Supply chain & Vendored dependencies at pinned revisions, license allow list enforced in continuous integration, signature verification on registry downloads \\
Credential handling & Provider credentials live in the platform keychain, are never written to the workspace, and are never included in context \\
Unbounded resource use & Every agent session and verification run carries a timeout, a token budget and a hard kill path \\
Data at rest & Workspace state is local. Nothing is uploaded without an explicit capability being enabled \\
\bottomrule
\caption{Security controls mapped to risks.}
\end{longtable}

\section{Privacy and telemetry}

Telemetry is off by default and is never a precondition for any feature. When a developer
enables it, the payload is enumerable in the settings surface before it is sent, contains no
source code, no file paths and no prompts, and can be inspected in a local log. Crash reports
are written locally and uploaded only on an explicit action.

\section{Accessibility}

Accessibility is an acceptance criterion on every interface issue rather than a milestone of
its own.

\begin{itemize}
  \item Every control has an accessible name and role, and every surface has a logical focus
  order.
  \item Every command is reachable from the keyboard, and the palette is the guaranteed path.
  \item Contrast meets the standard threshold in both themes, verified by an automated check
  in continuous integration.
  \item Severity is never conveyed by colour alone; it always carries a glyph and a word.
  \item Motion respects the platform reduced motion preference, including the activity pulse.
  \item Screen reader output is verified for the Mission Control and Verification Deck flows,
  because those are the surfaces that report autonomous work.
\end{itemize}

\section{Internationalisation}

Strings are externalised from the first commit, because retrofitting is expensive. The
initial shipping locales are English and Turkish. Locale affects interface strings only;
identifiers, log output and journal entries remain in English so that artefacts stay
portable across teams.

\section{Platform support}

\begin{table}[h]
\centering
\small
\begin{tabularx}{\textwidth}{l l X}
\toprule
\textbf{Platform} & \textbf{Status} & \textbf{Notes} \\
\midrule
macOS arm64 & Primary & Reference development and benchmark target \\
macOS x86\_64 & Supported & Built from the same source, verified in continuous integration \\
Linux x86\_64 & Supported & Distributed as a portable single file bundle \\
Linux arm64 & Best effort & Built, tested less frequently \\
Windows x86\_64 & Supported & Installer, path length and process semantics handled explicitly \\
\bottomrule
\end{tabularx}
\caption{Platform support matrix.}
\end{table}

\section{Packaging and updates}

\begin{itemize}
  \item Artefacts are produced by continuous integration from a clean checkout, never from a
  developer machine.
  \item Every artefact carries a checksum and, where the platform supports it, a signature.
  \item Updates are staged: a fraction of installations receive a release first, and a
  regression halts the rollout.
  \item An update never rewrites workspace state in place without a backup that the previous
  version can still read.
  \item An air gapped installation path exists, because enterprise customers require it and
  it also proves the build has no hidden network dependency.
\end{itemize}

\section{Risk register}

\begin{longtable}{@{}p{0.30\textwidth} p{0.14\textwidth} p{0.48\textwidth}@{}}
\toprule
\textbf{Risk} & \textbf{Severity} & \textbf{Mitigation} \\
\midrule
\endhead
Rendering framework churn upstream & Medium & Vendored and pinned, wrapped behind an internal crate so it can be replaced \\
Model provider terms change & High & Vendor routing is configurable per role, local models are a first class path, no feature depends on one provider \\
Autonomous run produces plausible but wrong work & High & Verification gate before display, adversarial review, one commit per ticket, journal, easy revert \\
Context engine underperforms naive stuffing & Medium & Permanent evaluation harness measuring first attempt success, with the naive strategy as the baseline \\
Scope creep into a general purpose editor & High & Non goals stated in Chapter one, every issue traceable to a story \\
Single maintainer bottleneck & High & Category partition designed for parallel work, documentation as the source of truth rather than tribal knowledge \\
Performance regression as features land & Medium & Budgets enforced by benchmark in continuous integration \\
Extension ecosystem security incident & Medium & Capability sandbox, signature verification, isolation on trap \\
License contamination through a transitive dependency & High & Allow list enforced mechanically on every build, vendored tree audited \\
\bottomrule
\caption{Risk register with mitigations that are implemented rather than aspirational.}
\end{longtable}
